%%%%%%%%%%%%%%%%%%%%%%%%%%%%%%%%%%%%%%%%%%%%%%%%%%%%%%%%%%%%%%%%%%%%%%
% Lớp 10 | Phương pháp tọa độ trong mặt phẳng (𝑂𝑥𝑦) | Đường thẳng trong mặt phẳng tọa độ
% Dạng 0H9?3-1: Điểm, véctơ, hệ  số góc của đường thẳng
% Số câu: 43  (sinh tự động bởi tach_cau.py; từ nay thầy sửa trực tiếp file này)
%%%%%%%%%%%%%%%%%%%%%%%%%%%%%%%%%%%%%%%%%%%%%%%%%%%%%%%%%%%%%%%%%%%%%%

%%%%%%%%%%%%%%%% TRẮC NGHIỆM 4 PHƯƠNG ÁN %%%%%%%%%%%%%%%%

% nguồn: MaHoaNH_TracNghiem_Lop10
\begin{ex}%[0H9N3-1]%[ma:2c8b388f]
	Một đường thẳng có  bao nhiêu vectơ chỉ phương?
	\choice 
		{$1$} 
		{$2$} 
		 {$3$}
		{\True Vô  số}
	\loigiai{} 
\end{ex}

% nguồn: MaHoaNH_TracNghiem_Lop10
\begin{ex}%[0H9N3-1]%[ma:d803cf31]
	Cho đường thẳng $\Delta\colon x-3y-2=0$. Tọa độ của vectơ nào sau đây \textbf{không phải} là tọa độ vectơ pháp tuyến của  $\Delta$?
	\choice 
		 {$(1;-3)$}
		{$(-2;6)$} 
		{$\left(  \dfrac{1}{3};-1\right) $}
		{\True $(3;1)$} 
	\loigiai{Ta có $\left(3;1\right)$ là tọa độ vectơ chỉ phương của $\Delta$} 
\end{ex}

% nguồn: MaHoaNH_TracNghiem_Lop10
\begin{ex}%[0H9N3-1]%[ma:64e2768b]
	Trong mặt phẳng tọa độ $Oxy$,  cho đường thẳng $d\colon x-2y+3=0$. Vectơ pháp tuyến của đường thẳng $d$ là
	\choice 
		 {\True $\overrightarrow{n}=(1;-2)$}
		 {$\overrightarrow{n}=(2; 1)$}
		{$\overrightarrow{n}=(-2; 3)$} 
		{$\overrightarrow{n}=(1;  3)$}
	 \loigiai{Vectơ pháp tuyến của đường thẳng $d\colon x-2y+3=0$ là $\overrightarrow{n}=(1;-2)$.}
\end{ex}

% nguồn: MaHoaNH_TracNghiem_Lop10
\begin{ex}%[0H9N3-1]%[ma:ebc36fdf]
	Trong mặt phẳng tọa độ $Oxy$,  cho đường thẳng $d\colon x-2y+2024=0$. Một vectơ pháp tuyến của đường thẳng $d$ có tọa độ là
	\choice 
		{\True $(-1;  2)$}
		{$(2;  1)$}
		{$(1;  2)$}
		{$(-2;  1)$}
	 \loigiai{
	Ta có một vectơ  pháp tuyến của đường thẳng $d\colon x-2y+2024=0$ là $(-1; 2)$.}
\end{ex}

% nguồn: MaHoaNH_TracNghiem_Lop10
\begin{ex}%[0H9N3-1]%[ma:93d6c9a2]
	 Phương trình tham số của đường thẳng $d$ đi qua điểm $M(-2;3)$ và có một vectơ chỉ phương $\overrightarrow{u}=(3;-4)$ là
	 \choice
		{$\heva{&x=-2+4t\\&y=3+3t}$} 
		{\True  $\heva{&x=-2+3t\\&y=3-4t}$}
		{$\heva{&x=-2+3t\\&y=3+4t}$} 
		 {$\heva{&x=5+4t\\&y=6-3t}$}
	\loigiai{Đường thẳng $d$ đi  qua điểm $M(-2;3)$ và có một vectơ chỉ phương $\overrightarrow{u}=(3;-4)$.\\
	Suy ra phương trình  tham số $d\colon\heva{&x=-2+3t\\&y=3-4t.}$}
\end{ex}

% nguồn: MaHoaNH_TracNghiem_Lop10
\begin{ex}%[0H9N3-1]%[ma:e4e346c2]
	Cho  $(d)\colon \heva{&{x=2+3t} \\&{y=5-4t}}$. Điểm nào sau đây là điểm không thuộc $(d)$
	\choice 
		{$A(5;3)$} 
		{\True $B(2;5)$} 
		{$C(-1;  9)$}
		{$D(8;-3)$} 
	\loigiai{ 
	Thay $B\left(2;5\right)\Rightarrow \heva{&{2=2+3t} \\&{5=5-4t}} \Rightarrow \heva{&{t=0} \\&{t=0}}  \Rightarrow t=0$.}
\end{ex}

% nguồn: MaHoaNH_TracNghiem_Lop10
\begin{ex}%[0H9N3-1]%[ma:361b76d1]
	Một vectơ pháp  tuyến của đường thẳng $2x-3y+6=0$ là
	\choice 
		 {\True $\overrightarrow{n_4}=\left(2;-3\right)$}
		 {$\overrightarrow{n_2}=\left(2; 3\right)$}
		 {$\overrightarrow{n_3}=\left(3; 2\right)$}
		{$\overrightarrow{n_1}=\left(-3;  2\right)$}
	\loigiai{ 
		Từ phương trình tổng quát  ta thấy một vectơ pháp tuyến của đường thẳng là $\overrightarrow{n_4}=\left(2;-3\right)$.
	} 
\end{ex}

% nguồn: MaHoaNH_TracNghiem_Lop10
\begin{ex}%[0H9N3-1]%[ma:29e20de6]
	Trong mặt phẳng tọa  độ $Oxy$, cho đường thẳng $\Delta\colon 3x+y-4=0$. Tọa độ một vectơ chỉ phương của đường thẳng $\Delta$ là
	 \choice
		{$\overrightarrow{u}_1=(3; 1)$} 
		 {\True $\overrightarrow{u}_2=(1;-3)$}
		{$\overrightarrow{u}_3=(3;-1)$} 
		 {$\overrightarrow{u}_4=(-1;-3)$}
	\loigiai{ 
		Đường thẳng $\Delta$ có một vectơ pháp tuyến là $\overrightarrow{n}=(3;  1)$ nên tọa độ của một vectơ chỉ phương của $\Delta$ là $\overrightarrow{u}_2=(1;-3)$.
	 }
\end{ex}

% nguồn: MaHoaNH_TracNghiem_Lop10
\begin{ex}%[0H9N3-1]%[ma:9e8132dd]
	Cho đường thẳng $\Delta\colon x+3y-4=0$. Điểm  nào dưới đây không thuộc $\Delta$?\\
	\choice 
		 {$M(1;1)$}
		 {$N(4;0)$}
		 {\True $P(2;-1)$}
		 {$Q(-2;2)$}
	 \loigiai{
		Thay lần lượt tọa độ các điểm vào phương trình đường thẳng $\Delta$ ta  thấy điểm $P$ không thuộc $\Delta$ vì $2+3\cdot(-1)-4=-5\neq 0$.\\
	} 
\end{ex}

% nguồn: MaHoaNH_TracNghiem_Lop10
\begin{ex}%[0H9N3-1]%[ma:9d752a95]
	vectơ nào dưới đây là một vectơ chỉ phương của đường  thẳng $\Delta\colon \heva{&{x=5-\dfrac{1}{2} t} \\ &{y=-3+3t}}$?
	\choice 
		{$\overrightarrow{u_1}=\left(-1;3\right)$} 
		 {$\overrightarrow{u_2}=\left(\dfrac{1}{2};3\right)$}
		{$\overrightarrow{u_3}=\left(5;-3\right)$} 
		 {\True $\overrightarrow{u_4}=\left(1;-6\right)$}
	\loigiai{ 
		Từ phương trình tham số ta thấy một vectơ chỉ phương của $\Delta\colon \heva{&{x=5-\dfrac{1}{2} t} \\&{y=-3+3t}}$ là $\overrightarrow{u}=\left(-\dfrac{1}{2};3\right)$  nên ta có thể chọn một vectơ chỉ phương là $\overrightarrow{u_4}=\left(1;-6\right)$.
	} 
\end{ex}

% nguồn: MaHoaNH_TracNghiem_Lop10
\begin{ex}%[0H9N3-1]%[ma:9a12d004]
	Trong mặt phẳng $Oxy$, cho đường thẳng $(d)\colon 5x-2y+8=0$. vectơ pháp tuyến của  đường thẳng $(d)$ là
	 \choice
		{$\vec{n}=(-2;-5)$} 
		 {$\vec{n}=(5; 2)$}
		{$\vec{n}=(2; 5)$} 
		{\True $\vec{n}=(5;-2)$} 
	\loigiai{ 
		Từ phương trình tổng quát ta có vectơ pháp tuyến  của đường thẳng $(d)$ là $\vec{n}=(5;-2)$.
	 }
\end{ex}

% nguồn: MaHoaNH_TracNghiem_Lop10
\begin{ex}%[0H9N3-1]%[ma:b3723abf]
	 Một vectơ chỉ phương của đường thẳng $\heva{&{x=2+3t} \\&{y=-3-t}}$ là
	\choice 
		{$\overrightarrow{u_1}=\left(2;-3\right)$} 
		{\True $\overrightarrow{u_2}=\left(3;-1\right)$} 
		{$\overrightarrow{u_3}=\left(3;  \; 1\right)$}
		{$\overrightarrow{u_4}=\left(3;-3\right)$} 
	 \loigiai{
		Từ phương trình tham số của đường thẳng ta có một vectơ  chỉ phương của đường thẳng là $\overrightarrow{u_2}=\left(3;-1\right)$.
	} 
\end{ex}

% nguồn: MaHoaNH_TracNghiem_Lop10
\begin{ex}%[0H9N3-1]%[ma:efa9aa36]
	Trong mặt phẳng  tọa độ $Oxy $, cho đường thẳng $d\colon \heva{&x=-4+5t\\ &y=3-t}\left( t\in\mathbb{R}\right) $. Vectơ nào sau đây là một vectơ chỉ phương của đường thẳng $d$?
	 \choice
		 {$\overrightarrow{n}_{3}=\left(5;1\right)$}
		 {\True $\overrightarrow{u}_{1}=\left(\dfrac{5}{2};-\dfrac{1}{2}\right)$}
		{$\overrightarrow{n}_{4}=\left(1;5\right)$} 
		{$\overrightarrow{u}_{2}=\left(-4;3\right)$} 
	 \loigiai{
		Từ phương trình tham số của $d$ ta có $\overrightarrow{u}=(5;-1)$ là một vectơ chỉ phương của $d$, ta thấy $\overrightarrow{u}_{1}=\left(\dfrac{5}{2};-\dfrac{1}{2}\right)=\dfrac{1}{2}\overrightarrow{u}  $. Do đó $\overrightarrow{u}_{1}=\left(\dfrac{5}{2};-\dfrac{1}{2}\right)$ cũng là vectơ chỉ phương của $d$.
	} 
\end{ex}

% nguồn: MaHoaNH_TracNghiem_Lop10
\begin{ex}%[0H9N3-1]%[ma:04fb9af1]
	Vectơ nào dưới đây là một vectơ chỉ phương của đường thẳng $d\colon \heva{&{x=2}  \\ &{y=-1+6t}}$?
	 \choice
		 {$\overrightarrow{u_1}=\left(6;0\right)$}
		 {$\overrightarrow{u_2}=\left(-6;0\right)$}
		{$\overrightarrow{u_3}=\left(2;6\right)$} 
		{\True $\overrightarrow{u_4}=\left(0;1\right)$} 
	 \loigiai{
		 Từ phương trình tham số ta thấy một vectơ chỉ phương của $d$ là $\overrightarrow{u}=\left(0;6\right)=6\left(0;1\right)$ nên ta có thể chọn một vectơ chỉ phương là $\overrightarrow{u_4}=\left(0;1\right)$.
	} 
\end{ex}

% nguồn: MaHoaNH_TracNghiem_Lop10
\begin{ex}%[0H9N3-1]%[ma:ed56a999]
	Đường thẳng $d\colon\heva{&x=-4+3t \\ 
	&y=1+2t}$ có  vectơ pháp tuyến có tọa độ là
	 \choice
		{$(1; 1)$} 
		{$(-4;-6)$} 
		{\True $(2;-3)$} 
		{$(-3;  2)$}
	 \loigiai{
	Đường thẳng $d$ có vectơ chỉ phương $\vec{u}=(3; 2)$ nên  vectơ pháp tuyến có tọa độ $(2;-3)$.}
\end{ex}

% nguồn: MaHoaNH_TracNghiem_Lop10
\begin{ex}%[0H9N3-1]%[ma:0b34ce14]
	Cho hai điểm $M(2; 3)$ và $N(-2; 5)$. Đường thẳng  $MN$ có một vectơ chỉ phương là
	\choice 
		{$\overrightarrow{u}=(4;  2)$}
		{\True $\overrightarrow{u}=(4;-2)$} 
		 {$\overrightarrow{u}=(-4;-2)$}
		{$\overrightarrow{u}=(-2;  4)$}
	 \loigiai{
		Ta có $\overrightarrow{MN}=(-4;  2)$. Do đó một vectơ chỉ phương của $MN$ là $\overrightarrow{u}=(4;-2)$.
	 }
\end{ex}

% nguồn: MaHoaNH_TracNghiem_Lop10
\begin{ex}%[0H9N3-1]%[ma:81382761]
	Cho đường thẳng $(d)\colon 3x-7y+15=0$. Mệnh đề nào sau  đây sai?
	\choice 
		{$\overrightarrow{u}=\left(7;3\right)$ là  vectơ chỉ phương của $(d)$}
		{$(d)$  có hệ số góc $k=\dfrac{3}{7}$}
		{$(d)$ không  đi qua góc tọa độ}
		{\True $(d)$ đi qua hai điểm $M\left(-\dfrac{1}{3};2\right)$ và $N\left(5;0\right)$} 
	 \loigiai{
		Giả sử $N\left(5;0\right)\in d\colon 3x-7y+15=0\Rightarrow 3\cdot 5-7\cdot 0+15=0$  (vô lí).
	} 
\end{ex}

% nguồn: MaHoaNH_TracNghiem_Lop10
\begin{ex}%[0H9N3-1]%[ma:0d6208e6]
	Cho đường thẳng $(d)\colon 2x+3y-4=0$. vectơ nào sau đây là vectơ pháp tuyến của $(d)$? 
	\choice 
		{$\overrightarrow{n_1}=\left(3;2\right)$} 
		 {\True $\overrightarrow{n_2}=\left(-4;-6\right)$}
		{$\overrightarrow{n_3}=\left(2;-3\right)$} 
		 {$\overrightarrow{n_4}=\left(-2;3\right)$}
	 \loigiai{
		Ta có  $(d)\colon 2x+3y-4=0\Rightarrow$ vectơ pháp tuyến $\overrightarrow{n}=\left(2;3\right)=\left(-4;-6\right)$.
	 }
\end{ex}

% nguồn: MaHoaNH_TracNghiem_Lop10
\begin{ex}%[0H9N3-1]%[ma:46257ea6]
	Trong mặt phẳng $Oxy$, cho đường  thẳng $\Delta$ có phương trình $\dfrac{x-1}{2}=\dfrac{y+1}{-1}$. Vectơ chỉ phương của đường thẳng $\Delta$ là
	 \choice
		{\True $\overrightarrow{u}=(2;-1)$} 
		{$\overrightarrow{u}=(1;  2)$}
		 {$\overrightarrow{u}=(1;-1)$}
		 {$\overrightarrow{u}=(1; 1)$}
	 \loigiai{
		vectơ chỉ phương của  đường thẳng $\Delta$ là $\overrightarrow{u}=(2;-1)$.
	 }
\end{ex}

% nguồn: MaHoaNH_TracNghiem_Lop10
\begin{ex}%[0H9N3-1]%[ma:65021974]
	Cho phương trình $ax+by+c=0\; (1)$ với $a^2+b^2> 0$. Mệnh đề nào sau đây  sai?
	\choice 
		{$(1)$ là phương trình tổng quát của đường  thẳng có vectơ pháp tuyến là $\overrightarrow{n}=\left(a;b\right)$}
		{$a=0$ $(1)$ là  phương trình đường thẳng song song hoặc trùng với trục $Ox$}
		{$b=0$ $(1)$ là phương trình đường thẳng song song hoặc trùng với trục  $Oy$}
		{\True Điểm $M_0 \left(x_0;y_0 \right)$ thuộc đường thẳng $(1)$ khi  và chỉ khi $ax_0+by_0+c\ne 0$}
	 \loigiai{
		Ta có điểm $M_0 \left(x_0;y_0 \right)$ thuộc đường thẳng $(1)$ khi và chỉ khi  $ax_0+by_0+c=0$.
	} 
\end{ex}

% nguồn: MaHoaNH_TracNghiem_Lop10
\begin{ex}%[0H9N3-1]%[ma:541b3f31]
	Mệnh đề nào sau đây sai? Đường thẳng $(d)$ được xác định khi  biết.
	 \choice
		{\True  Một vectơ pháp tuyến hoặc một vectơ chỉ phương}
		{Hệ  số góc và một điểm thuộc đường thẳng}
		{Một điểm thuộc $(d)$ và biết $(d)$  song song với một đường thẳng cho trước}
		{Hai  điểm phân biệt thuộc $(d)$}
	 \loigiai{
		Nếu chỉ có vectơ pháp tuyến hoặc một vectơ chỉ  phương thì thiếu điểm đi qua để viết đường thẳng.
	} 
\end{ex}

% nguồn: MaHoaNH_TracNghiem_Lop10
\begin{ex}%[0H9N3-1]%[ma:8591f540]
	Trong mặt phẳng  tọa độ $Oxy$, cho đường thẳng $\Delta \colon 2x+5y+1=0$. Một véc-tơ pháp tuyến của đường thẳng $\Delta$ có tọa độ là
	\choice 
		 {$(5;1)$}
		 {\True $(2;5)$}
		 {$(-2;5)$}
		 {$(5;-2)$}
	\loigiai{ 
		Một véc-tơ pháp tuyến của đưòng thẳng $\Delta$  có tọa độ là $(2;5)$.
	 }
\end{ex}

% nguồn: MaHoaNH_TracNghiem_Lop10
\begin{ex}%[0H9N3-1]%[ma:1e3afba7]
	Trong mặt phẳng $Oxy$, cho đường thẳng $d\colon  \heva{&x=1-t\\&y=3+2t}$. Véc-tơ nào sau đây là một véc-tơ pháp tuyến của $d$?
	\choice 
		{$\overrightarrow{n}=(-1;2)$} 
		{$\overrightarrow{n}=(1;2)$} 
		{\True $\overrightarrow{n}=(2;1)$} 
		{$\overrightarrow{n}=(2;-1)$} 
	\loigiai{ 
		Đường thẳng $d$ có một véc-tơ chỉ phương là $\overrightarrow{a}=(-1;2)$ nên nhận một véc-tơ pháp tuyến là  $\overrightarrow{n}=(2;1)$.
	} 
\end{ex}

% nguồn: MaHoaNH_TracNghiem_Lop10
\begin{ex}%[0H9N3-1]%[ma:e2ee6cf0]
	Trong mặt phẳng toạ độ $Oxy$, cho đường thẳng $d\colon  2x+3y+10=0$ có véc-tơ pháp tuyến là
	\choice 
		{\True  $\overrightarrow{n}=\left( 2;3 \right)$}
		{$\overrightarrow{n}=\left(  3;2 \right)$}
		 {$\overrightarrow{n}=\left( 2;10 \right)$}
		 {$\overrightarrow{n}=\left( 3;10 \right)$}
	 \loigiai{
		Đường thẳng  $d$ có véc-tơ pháp tuyến là $\overrightarrow{n}=(2;3)$.
	 }
\end{ex}

% nguồn: MaHoaNH_TracNghiem_Lop10
\begin{ex}%[0H9N3-1]%[ma:6d645e18]
	Vec-tơ nào dưới đây là một vec-tơ chỉ  phương của đường thẳng $\Delta\colon\heva{&x=5-\dfrac{1}{2}t\\&y=-3+3t}$?
	 \choice
		{$\overrightarrow{u}_1=(5;-3)$} 
		 {$\overrightarrow{u}_2=\left(\dfrac{1}{2};3\right)$}
		{\True  $\overrightarrow{u}_3=\left(-\dfrac{1}{2};3\right)$}
		 {$\overrightarrow{u}_4=(-1;-6)$}
	\loigiai{ 
		 Một vec-tơ chỉ phương của đường thẳng là $\overrightarrow{u}_3=\left(-\dfrac{1}{2};3\right)$.
	} 
\end{ex}

% nguồn: MaHoaNH_TracNghiem_Lop10
\begin{ex}%[0H9N3-1]%[ma:dfe17d45]
	 vectơ chỉ phương của đường thẳng $\dfrac{x}{3}+\dfrac{y}{2}=1$ là
	\choice 
		{$\overrightarrow{u}_4=\left(-2;  3\right)$}
		{\True $\overrightarrow{u}_2=\left(3;-2\right)$} 
		{$\overrightarrow{u}_3=\left(3; 2\right)$} 
		 {$\overrightarrow{u}_1=\left(2; 3\right)$}
	\loigiai{ 
		Ta có $\dfrac{x}{3}+\dfrac{y}{2}=1\Leftrightarrow 2x+3y-6=0$ nên đường thẳng có vectơ pháp tuyến là  $\overrightarrow{n}=\left(2; 3\right)$.\\
		 Suy ra vectơ chỉ phương là $\overrightarrow{u}=\left(3;-2\right)$.
	} 
\end{ex}

% nguồn: MaHoaNH_TracNghiem_Lop10
\begin{ex}%[0H9N3-1]%[ma:4cf3ef82]
	Cho đường thẳng $\Delta$ có phương trình tổng quát $-2x+3y-1=0$. vectơ nào  sau đây không là vectơ chỉ phướng của $\Delta$.
	 \choice
		 {$\left(1; \dfrac{2}{3}\right)$}
		{$(3;  2)$}
		{\True  $(2; 3)$}
		 {$(-3;-2)$}
	\loigiai{ 
		Từ phương trình tổng quát ta thấy một vectơ pháp tuyến là  $\overrightarrow{n}=(-2;3)$ suy ra một vectơ chỉ phương của đường thẳng là $\overrightarrow{u}=(3;2)=-1(-3;-2)=3\left(1; \dfrac{2}{3}\right)$.\\
		Vậy vectơ có  tọa độ $(2;3)$ không phải là vectơ chỉ phương của $\Delta$.
	 }
\end{ex}

% nguồn: MaHoaNH_TracNghiem_Lop10
\begin{ex}%[0H9N3-1]%[ma:c718dfb2]
	Cho đường thẳng  $\Delta$ có phường trình tổng quát $-2x+3y-1=0$.
	vectơ nào sau đây là vectơ chỉ phương của đường  thẳng $\Delta$.
	 \choice
		{\True $(3; 2)$} 
		 {$(2; 3)$}
		{$(-3; 2)$} 
		 {$(2;-3)$}
	 \loigiai{
		Từ phương trình tổng quát ta thấy một vectơ pháp tuyến của $\Delta$ là $\overrightarrow{n}=(-2; 3)$ suy ra một vectơ chỉ  phương là $\overrightarrow{u}=(3;2)$.
	 }
\end{ex}

% nguồn: MaHoaNH_TracNghiem_Lop10
\begin{ex}%[0H9N3-1]%[ma:81004e64]
	Cho đường thẳng $\Delta\colon \heva{&x=1+2t\\&y=3-t}\,(t\in \mathbb{R})$. Vectơ nào dưới đây là vectơ chỉ phương của đường  thẳng $\Delta$?\\
	\choice 
		{$\overrightarrow{u}_{1}=(1;3)$} 
		 {$\overrightarrow{u}_{2}=(2;0)$}
		{$\overrightarrow{u}_{3}=(2;1)$} 
		 {\True $\overrightarrow{u}_{4}=(2;-1)$}
	\loigiai{ 
		Ta có  $\Delta\colon \heva{&x=1+2t\\&y=3-t}\,(t\in \mathbb{R})$.\\
		Vectơ chỉ phương của $\Delta$  là $\overrightarrow{u}_{4}=(2;-1)$.
	 }
\end{ex}

% nguồn: MaHoaNH_TracNghiem_Lop10
\begin{ex}%[0H9N3-1]%[ma:ce6cfa9d]
	Đường thẳng  $(d)$ có vectơ pháp tuyến $\overrightarrow{n}=\left(a;b\right)$. Mệnh đề nào sau đây sai?
	\choice 
		{$\overrightarrow{u}_1=\left(b;-a\right)$ là vectơ chỉ phương của $(d)$} 
		{$\overrightarrow{u}_2=\left(-b;a\right)$ là vectơ  chỉ phương của $(d)$}
		{$\overrightarrow{n'}=\left(ka;kb\right) k\in R$ là  vectơ pháp tuyến của $(d)$}
		{\True $(d)$ có  hệ số góc $k=\dfrac{-b}{a} \left(b\ne 0\right)$}
	\loigiai{ 
		 Phương trình đường thẳng có vectơ pháp tuyến $\overrightarrow{n}=\left(a;b\right)$ là $ax+by+c=0\Leftrightarrow y=-\dfrac{a}{b} x-\dfrac{c}{b} \left(b\ne 0\right)$.\\
		 Suy ra hệ số góc $k=-\dfrac{a}{b}$.
	 }
\end{ex}

% nguồn: MaHoaNH_TracNghiem_Lop10
\begin{ex}%[0H9N3-1]%[ma:2b58e6d5]
	Trong mặt phẳng tọa độ $Oxy$, cho tam  giác $ABC$ có trọng tâm $G$. Tìm tọa độ điểm $C$ biết $A(3;1)$, $B(-4;0)$, $G(1;1)$?
	 \choice
		{$C(2;0)$} 
		 {$C(-2;0)$}
		 {$C(4;-2)$}
		{\True  $C(4;2)$}
	 \loigiai{
		Vì $G$ là trọng tâm  của tam giác nên ta có $$\heva{&x_G=\dfrac{x_A+x_B+x_C}{3}\\ &y_G=\dfrac{y_A+y_B+y_C}{3}} \Rightarrow \heva{&x_C=3x_G-x_A-x_B=3 \cdot 1 -3-(-4)=4\\&y_C=3y_G-y_A-y_B=3 \cdot 1 -1-0=2.}$$
		Vậy  $C(4;2)$.
	} 
\end{ex}

% nguồn: MaHoaNH_TracNghiem_Lop10
\begin{ex}%[0H9N3-1]%[ma:a388261c]
	Cho đường thẳng $(d)\colon \heva{&{x=2-3t} \\&{y=-1+2t}}$ và  điểm $A\left(\dfrac{7}{2};-2\right)$. Điểm $A\in (d)$ ứng với giá trị nào
	 \choice
		 {$t=\dfrac{3}{2}$}
		 {$t=\dfrac{1}{2}$}
		 {\True $t=-\dfrac{1}{2}$}
		 {$t=2$}
	 \loigiai{
		Ta có $A\left(\dfrac{7}{2};-2\right)\in  (d)\Rightarrow \heva{&{\dfrac{7}{2}=2-3t} \\
		&{-2=-1+2t} 
		}  \Rightarrow \heva{&{t=-\dfrac{1}{2}} \\
		 &{t=-\dfrac{1}{2}}
		}  \Rightarrow t=-\dfrac{1}{2}$.
	 }
\end{ex}

% nguồn: MaHoaNH_TracNghiem_Lop10
\begin{ex}%[0H9H3-1]%[ma:b115d38b]
	Trong mặt phẳng tọa độ $Oxy$, cho hai điểm $A(1;-2)$ và $B(3; 2)$. Phương trình tổng quát của đường thẳng  $AB$ là
	 \choice
		{$2x+4y+6=0$} 
		{$2x-y+4=0$} 
		{$x+2y-10=0$} 
		 {\True $2x-y-4=0$}
	\loigiai{ 
		Ta có đường thẳng $AB$ nhận $\overrightarrow{AB}=(2; 4)$ là một vectơ chỉ phương nên đường thẳng $AB$ có  một vectơ pháp tuyến là $\overrightarrow{n}=(2;-1)$.\\
		Phương trình tổng quát  của đường thẳng $AB$ là\\ \centerline{$2(x-1)-1(y+2)=0 \Leftrightarrow 2x-y-4=0$.}
	} 
\end{ex}

% nguồn: MaHoaNH_TracNghiem_Lop10
\begin{ex}%[0H9H3-1]%[ma:08a15a03]
	vectơ nào dưới đây là một vectơ chỉ phương  của đường thẳng đi qua hai điểm $A\left(-3;2\right)$ và $B\left(1;4\right)$?
	\choice 
		{$\overrightarrow{u_1}=\left(-1;2\right)$} 
		 {\True $\overrightarrow{u_2}=\left(2;1\right)$}
		 {$\overrightarrow{u_3}=\left(-2;6\right)$}
		{$\overrightarrow{u_4}=\left(1;1\right)$} 
	\loigiai{ 
		Ta có $\overrightarrow{AB}=\left(4;2\right)$ một vectơ chỉ phương  của đường thẳng $AB$ cùng phương với $\overrightarrow{AB}=\left(4;2\right)$.\\
		Ta thấy  $\overrightarrow{u_2}=\left(2;1\right)=\dfrac{1}{2} \overrightarrow{AB}$.\\
		Vậy $\overrightarrow{u_2}=\left(2;1\right)$ là một  vectơ chỉ phương của $AB$.
	} 
\end{ex}

% nguồn: MaHoaNH_TracNghiem_Lop10
\begin{ex}%[0H9H3-1]%[ma:386e543e]
	vectơ nào dưới đây là một vectơ  pháp tuyến của đường thẳng đi qua hai điểm $A\left(2;3\right)$ và $B\left(4;1\right)$?
	\choice 
		 {$\overrightarrow{n_1}=\left(2;-2\right)$}
		{$\overrightarrow{n_2}=\left(2;-1\right)$} 
		{\True  $\overrightarrow{n_3}=\left(1;1\right)$}
		 {$\overrightarrow{n_4}=\left(1;-2\right)$}
	 \loigiai{
		Ta có $\overrightarrow{AB}=\left(2;-2\right)$ một vectơ pháp  tuyến $\overrightarrow{n}$ của đường thẳng $AB$ thì vuông góc với $AB$.\\
		Suy ra $\overrightarrow{n}\cdot  \overrightarrow{AB}=0\Leftrightarrow x\cdot 2+y\cdot \left(-2\right)=0$.\\
		Chọn  $x=1,y=1\Rightarrow \overrightarrow{n}=\left(1;1\right)$.
	 }
\end{ex}

% nguồn: MaHoaNH_TracNghiem_Lop10
\begin{ex}%[0H9Y3-1]%[ma:7000214a]
	Cho đường tròn $(C)$ có phương trình  $(x+3)^2 + (y-4)^2 =2$. Tìm tâm và bán kính của đường tròn $(C)$.
	\choice 
		 {\True $I(-3; 4), R=\sqrt{2}$}
		 {$I(-3; -4), R=2$}
		 {$I(3; -4), R=\sqrt{2}$}
		{$I(3; -4),  R=2$}
	\loigiai{ 
		Từ  phương trình đường tròn $(C)$, ta suy ra đường tròn $(C)$ có tâm $I(-3;4)$ và bán kính $R=\sqrt{2}$.
	} 
\end{ex}

% nguồn: MaHoaNH_TracNghiem_Lop10
\begin{ex}%[0H9Y3-1]%[ma:ed2d3194]
	Đường thẳng $\Delta$ có véc-tơ chỉ phương là $\overrightarrow{u}_{\Delta}=(12;-13)$. Véc-tơ nào sau đây là véc-tơ  pháp tuyến của đường thẳng $\Delta$?
	\choice 
		 {$\overrightarrow{n}_{\Delta}=(-13;12)$}
		{$\overrightarrow{n}_{\Delta}=(12;13)$} 
		 {\True $\overrightarrow{n}_{\Delta}=(13;12)$}
		 {$\overrightarrow{n}_{\Delta}=(-12;-13)$}
	\loigiai{ 
		Véc-tơ pháp  tuyến của đường thẳng $\Delta$ là $\overrightarrow{n}_{\Delta}=(13;12)$.
	} 
\end{ex}


%%%%%%%%%%%%%%%% ĐÚNG/SAI %%%%%%%%%%%%%%%%

% nguồn: MaHoaNH_DungSai_Lop10
\begin{ex}%[0H9N3-1]%[ma:79e61a37]
	Trong hệ tọa độ  $Oxy$, các khẳng định sau đây đúng hay sai?
	 \choiceTF
		{\True  Hệ số góc của đường thẳng $y=2x-1$ là $k=2$}
		{\True Một véc-tơ pháp tuyến của đường thẳng song song với trục  $Ox$ là $\overrightarrow{u}=(0;1)$}
		{Điểm  $A(2;1)$ nằm trên đường thẳng $\Delta\colon 5x+3y-7=0$}
		{Một véc-tơ chỉ phương của đường thẳng $d\colon\heva{& x=1-4t\\& y=-2+3t}$ $(t\in\mathbb{R})$  là $\overrightarrow{u}=(1;-2)$}
	 \loigiai{
		\begin{itemchoice}
			\itemch Đường thẳng $y=2x-1$ có  hệ số góc là $k=2$.
			\itemch Các đường thẳng song song với trục $Ox$ có một véc-tơ  pháp tuyến là $\overrightarrow{u}=\overrightarrow{j}=(0;1)$.
			\itemch Thay tọa độ điểm $A$ vào  phương trình $\Delta$, ta có
			$$5\cdot 2+3\cdot  1-7=6\neq 0 \Rightarrow A\notin\Delta.$$
			\itemch Một véc-tơ chỉ phương của đường thẳng  $d\colon\heva{& x=1-4t\\& y=-2+3t}$ $(t\in\mathbb{R})$ là $\overrightarrow{a}=(-4;3)$.\\
			Do $1\cdot 3\neq (-2)\cdot (-4)$ nên $\overrightarrow{u}$ không cùng phương với  $\overrightarrow{a}$.\\
			Suy ra $\overrightarrow{u}=(1;-2)$ không phải là một  véc-tơ chỉ phương của $d$.
		\end{itemchoice}
	 }
\end{ex}

% nguồn: MaHoaNH_DungSai_Lop10
\begin{ex}%[0H9N3-1]%[0H9N3-5]%[ma:a6ff2ea1]
	Trong mặt phẳng tọa độ $Oxy$, cho đường thẳng $\Delta \colon x-2y+3=0$ và điểm $M(-2;1)$. 
	\choiceTF 
		{\True Đường thẳng $\Delta$ có một vectơ pháp tuyến là $\overrightarrow{u}=(3;-6)$} 
		{\True Đường  thẳng $\Delta$ có một vectơ chỉ phương là $\overrightarrow{n}=(2;1)$}
		{\True Khoảng  cách từ điểm $M$ đến đường thẳng $\Delta$ bằng $\dfrac{\sqrt{5}}{5}$}
		{Giao điểm của đường thẳng $\Delta$ với  trục $Oy$ là điểm $B(3;0)$}
	\loigiai{ 
		\begin{itemchoice}
			\itemch Ta có $\Delta \colon x-2y+3=0$  suy ra $\Delta$ có một vectơ pháp tuyến là $\overrightarrow{n}_{\Delta}=(1;-2)$.\\
			Ta lại có $\overrightarrow{u}=(3;-6)=3(1;-2)=3\overrightarrow{n}_{\Delta}$.\\ 
			 Suy ra $\overrightarrow{n}_{\Delta}$ và $\overrightarrow{u}$ cùng phương.\\
			 Suy ra $\overrightarrow{u}$ cũng là một vectơ pháp tuyến của đường thẳng $\Delta$.
			\itemch Do $\Delta$ có một vectơ pháp tuyến là $\overrightarrow{n}_{\Delta}=(1;-2)$ nên đường  thẳng $\Delta$ có một vectơ chỉ phương là $\overrightarrow{n}=(2;1)$.
			\itemch Khoảng  cách từ điểm $M$ đến đường thẳng $\Delta$ là
			 \[\mathrm{d}(M,\Delta)=\dfrac{|1\cdot (-2)-2\cdot 1+3|}{\sqrt{1^2+(-2)^2}}=\dfrac{\sqrt 5}{5}.\]
			\itemch Ta  có thay $x=0$ vào phương trình của $\Delta \colon x-2y+3=0$ ta được $y=3$.\\
			Vậy giao điểm của đường thẳng $\Delta$  với trục $Oy$ là điểm $(0;3)$.
		\end{itemchoice}
	} 
\end{ex}

% nguồn: MaHoaNH_DungSai_Lop10
\begin{ex}%[0H9N3-1]%[0H9H3-5]%[0H9N3-6]%[ma:117d7b06]
	Trong mặt  phẳng tọa độ $Oxy$, cho đường thẳng $ \Delta\colon 3x+4y-1=0 $. Xét tính Đ/S các khẳng định sau:
	\choiceTF 
		{\True $\Delta$ có  một véc-tơ pháp tuyến là $\overrightarrow{n}=(6;8)$}
		{ \True Khoảng cách  từ $M(-4;3)$ đến $\Delta$ bằng $\dfrac{1}{5}$}
		{$\Delta$ đi qua  $M(1;1)$}
		{$\Delta$ tạo với trục hoành một góc bằng $30^\circ$} 
	 \loigiai{
		\begin{itemize}
			\item Đường thẳng $ \Delta\colon 3x+4y-1=0 $ có một véc-tơ pháp tuyến là $\overrightarrow{n_1}=(3;4)$ nên $\overrightarrow{n}=(6;8)$ cũng là một véc-tơ pháp tuyến của $\Delta$ vì  $\overrightarrow{n}=(6;8)$ cùng phương với $\overrightarrow{n_1}$.
			\item Khoảng  cách từ $M(-4;3)$ đến $\Delta$ là
			 $$\mathrm{d}(M,\Delta)=\dfrac{|3\cdot (-4)+4\cdot 3-1|}{\sqrt{(-4)^2+3^2}}=\dfrac{1}{5}.$$
			\item $M(1;1)$ không thuộc  $\Delta$ vì $3\cdot 1+4\cot 1-1=6\ne 0$.
			\item Ta có hệ số góc của $\Delta$ là $k=\dfrac{-3}{4}$ nên góc tạo bởi $\Delta$ và $Ox$ không bằng $30^\circ$ vì $\tan  30^\circ=\dfrac{1}{\sqrt{3}}$.
		\end{itemize}
		\begin{itemchoice}
			 \itemch Đúng.
			\itemch  Đúng.
			\itemch  Sai.
			\itemch Sai. 
		\end{itemchoice}
	} 
\end{ex}

% nguồn: MaHoaNH_DungSai_Lop10
\begin{ex}%[0H9H3-1]%[ma:749466fd]
	Trong hệ tọa độ $Oxy$, cho đường thẳng $d$ có một véc-tơ pháp tuyến là $\overrightarrow{n}=(a; b)$, với $a$, $b\in\mathbb{R}$. Trong các  khẳng định sau, khẳng định nào đúng, khẳng định nào sai?
	\choiceTF 
		 {\True Nếu $b=0$ thì đường thẳng $d$ không có hệ số góc}
		{Nếu $b\neq 0$ thì hệ số góc của đường thẳng $d$ là  $k=\dfrac{a}{b}$}
		{\True Đường thẳng $d$ có một véc-tơ chỉ  phương là $\overrightarrow{u}=(b;-a)$}
		{Véc-tơ $k\overrightarrow{n}$, $(k\in\mathbb{R})$ là  véc-tơ pháp tuyến của $d$}
	 \loigiai{
		$d$ có một véc-tơ  pháp tuyến là $\overrightarrow{n}=(a; b) \Rightarrow$ phương trình $d\colon ax+by+c=0$.
		\begin{itemize}
			\item Nếu $b=0$ thì đường thẳng  $d\colon ax+c=0$ không có hệ số góc.
			\item Nếu $b\neq 0$ thì đường thẳng $d\colon y=-\dfrac{a}{b} x-\dfrac{c}{b}$ có hệ số góc là $k=-\dfrac{a}{b}$. 
			\item Với $\overrightarrow{u}=(b;-a)  \Rightarrow \overrightarrow{u}\cdot\overrightarrow{n}=0 \Rightarrow \overrightarrow{u}\perp\overrightarrow{n}$\\
			 $\Rightarrow\overrightarrow{u}$ là một véc-tơ chỉ phương của $d$.
			\item  Chọn $k=0 \Rightarrow k\overrightarrow{n}=(0; 0)$ không phải là véctơ pháp tuyến của $d$.
		\end{itemize}
	} 
\end{ex}

% nguồn: MaHoaNH_DungSai_Lop10
\begin{ex}%[0H9V3-1]%[ma:8e8ba9c8]
	Trong hệ tọa độ $Oxy$, cho tam  giác $ABC$ có $A(-1;2)$, $B(3;2)$, $C(4;6)$
	\choiceTF{\True Trọng tâm $G$ của  tam giác $ABC$ có tọa độ là $\left(2;\dfrac{10}{3}\right)$}
		{Trung điểm $M$ của đoạn $AB$ có  tọa độ là $(1;-2)$}
		{Đường  cao $AH$ của tam giác $ABC$ có phương trình là $x+4y+7=0$}
		{\True  Tâm đường tròn ngoại tiếp tam giác $ABC$ có tọa độ là $\left(1;\dfrac{37}{8}\right)$}
	\loigiai{ 
		\begin{itemchoice}
			\itemch Vì $G$ là trọng tâm  của tam giác $ABC$ nên ta có
			$$\heva{&x_G=\dfrac{-1+3+4}{3}=2\\&y_G=\dfrac{2+2+6}{3}=\dfrac{10}{3}.}$$ 
			 Vậy $G\left(2;\dfrac{10}{3}\right)$.
			\itemch  Vì $M$ là trung điểm của đoạn $AB$ nên ta có
			$$\heva{&x_M=\dfrac{-1+3}{2}=1\\&y_M=\dfrac{2+2}{2}=2.}$$ 
			 Vậy $M\left(1;2\right)$.
			\itemch Vì  $AH\perp BC$ nên $AH$ có vectơ pháp tuyến là $\overrightarrow{BC}=(1;4)$.\\
			Phương  trình tổng quát của $AH$ là $1(x+1)+4(y-2)=0\Leftrightarrow x+4y-7=0$.
			\itemch Gọi $(C)$ là đường tròn ngoại tiếp tam giác  $ABC$.\\
			Phương trình $(C)$ có dạng $x^2+y^2-2ax-2by+c=0$  với $a^2+b^2-c>0$.\\
			Vì $(C)$ đi qua $3$ điểm $A$, $B$, $C$ nên ta  có hệ phương trình
			$$\heva{&(-1)^2+2^2-2\cdot a\cdot (-1)-2\cdot b\cdot 2+c=0\\&3^2+2^2-2\cdot a\cdot 3-2\cdot b\cdot 2+c=0\\&4^2+6^2-2\cdot a\cdot 4-2\cdot  b\cdot 6+c=0}\Leftrightarrow \heva{&2a-4b+c=-5\\&-6a-4b+c=-13\\&-8a-12b+c=-52}\Leftrightarrow\heva{&a=1\\&b=\dfrac{37}{8}\\&c=\dfrac{23}{2}.}$$
			Vậy  tâm đường tròn ngoại tiếp tam giác $ABC$ có tọa độ là $\left(1;\dfrac{37}{8}\right)$.
		\end{itemchoice}
	 }
\end{ex}


%%%%%%%%%%%%%%%% TỰ LUẬN %%%%%%%%%%%%%%%%

% nguồn: MaHoaNH_TuLuan_Lop10
\begin{bt}%[0H9Y3-1]%[0H9Y3-3]%[ma:77a619ae]
	Trong mặt phẳng tọa độ $Oxy$, cho đường tròn $(C) \colon x^2+y^2-2x+6y-15=0$. 
	\begin{enumerate}
		\item Tìm tọa độ tâm $I$  và tính bán kính của đường tròn $(C)$.
		\item Viết phương trình tiếp tuyến $\Delta$ của đường tròn $(C)$ tại $A(4;1)$. 
	\end{enumerate}
	 \loigiai{
		\begin{enumerate}
			\item  Tâm $I(1;-3)$ và bán kính $R=\sqrt{1^2+(-3)^2-(-15)}=5$.
			\item Ta có $\overrightarrow{IA}=(3;4)$.\\ 
			Tiếp tuyến $\Delta$ đi qua  $A(4;1)$ nhận $\overrightarrow{IA}=(3;4)$ làm VTPT nên có phương trình là
			 $$3(x-4)+4(y-1)=0 \Leftrightarrow 3x+4y-16=0.$$
		\end{enumerate}
	} 
\end{bt}
